\documentclass{article}
\usepackage{amssymb}
\usepackage{amsfonts}
\usepackage{amsmath}

\author{Jonathan Kalsky}
\date{April 17 2026}
\title{\normalsize CSCE 222\\ \Large Homework 5}
\linespread{1.5}

\begin{document}
\maketitle
\tableofcontents
\vspace{1cm}

\subsection[12]{12 - Show that the sequence $\{a_n\}$ is a solution of the recurrence relation $a_n = -3a_{n-1} + 4a_{n-2}$ if}
P(n) = $a_n = -3{a_n-1} + 4a_{n-2}$ for all n$\geq$2\\
\subsubsection[a]{a - $a_n = 0$}
Basis step:
$P(2) = a_2 = 0 = -3(0) + 4(0) = 0$ which holds\\
Inductive step:
$p(k) \to p(k+1)$\\
Assume p(k) holds for k $\geq$ 2, $a_k = 0, a_{k-1} = 0$\\
LHS: $a_{k+1} = 0$\\
RHS: $-3{a_k} + 4a_{k-1}$\\
$-3{(0)} + 4{(0)} = 0$\\
LHS = RHS therefore $p(k) \to p(k+1)$ holds for all n $\geq$ 2 by induction.

\subsubsection[b]{b - $a_n = 1$}
Basis step:
$P(2) = a_2 = 1 = -3(1) + 4(1) = 1$ which holds\\
Inductive step:
$p(k) \to p(k+1)$\\
Assume p(k) holds for k $\geq$ 2, $a_k = 1, a_{k-1} = 1$\\
LHS: $a_{k+1} = 1$\\
RHS: $-3{a_k} + 4a_{k-1}$\\
$-3{(1)} + 4{(1)} = 1$\\
LHS = RHS therefore $p(k) \to p(k+1)$ holds for all n $\geq$ 2 by induction.

\subsubsection[c]{c - $a_n = (-4)^n$}

Basis step:
$P(2) = a_2 = (-4)^2 = 16 = -3((-4)^1) + 4((-4)^0) = 12 + 4 = 16$ which holds\\
Inductive step:
$p(k) \to p(k+1)$\\
Assume p(k) holds for k $\geq$ 2, $a_k = -4^k, a_{k-1} = -4^{k-1}$\\
LHS: $a_{k+1} = (-4)^{k+1}$\\
RHS: $-3{a_k} + 4a_{k-1}$\\
$-3{((-4)^k)} + 4{((-4)^{k-1})}$\\
$-3{((-4)^k)} + 4{((-4)^{k}(-4)^{-1})}$\\
$-3{((-4)^k)} + {((-4)^{k}(-1))}$\\
$-4{(-(4)^k)}$\\
$(-4)^{1}{((-4)^k)}$\\
${((-4)^{k+1})}$\\
LHS = RHS therefore $p(k) \to p(k+1)$ holds for all n $\geq$ 2 by induction.

\subsubsection[d]{d - $a_n = 2(-4)^n + 3$}
Basis step:
$P(2) = a_2 = 2(-4)^2 + 3 = 35 = -3(2(-4)^1 + 3) + 4(2(-4)^0 + 3) = 15 + 20 = 35$ which holds\\
Inductive step:
$p(k) \to p(k+1)$\\
Assume p(k) holds for k $\geq$ 2, $a_k = 2(-4)^k + 3, a_{k-1} = 2(-4)^{k-1} + 3$\\
LHS: $2(-4)^{k+1} + 3$\\
RHS: $-3{a_k} + 4a_{k-1}$\\
$-3{(2(-4)^k + 3)} + 4{(2(-4)^{k-1} + 3)}$\\
${(-6(-4)^k + -9)} + {(8(-4)^{k}(-4)^{-1} + 12)}$\\
${(-6(-4)^k + -9)} + {((-4)^{k}(-2) + 12)}$\\
${(-6(-4)^k)} + {(-2(-4)^{k}) + 3}$\\
${(-8(-4)^k)} + 3$\\
${((2(-4)^1)(-4)^k)} + 3$\\
${2(-4)^{k+1}} + 3$\\
LHS = RHS therefore $p(k) \to p(k+1)$ holds for all n $\geq$ 2 by induction.

\subsection{Question 14 - Find the recurrence relation satisfied by the sequence \{$a_n$\}}
\subsubsection[a]{a - $a_n = 0$}
$a_n$ = 0($a_{n-1}$) for all n $\geq$ 1, where $a_0$ = 0\\

\subsubsection[b]{b - $a_n = 1$}
$a_n$ = ($a_{n-1}$) for all n $\geq$ 1, where $a_0$ = 1\\

\subsubsection[c]{c - $a_n = 2^n$}
$a_n$ = 2($a_{n-1}$) for all n $\geq$ 1, where $a_0$ = 1\\

\subsubsection[d]{d - $a_n = 4^n$}
$a_n$ = 4($a_{n-1}$) for all n $\geq$ 1, where $a_0$ = 1\\

\subsubsection[e]{e - $a_n = n4^n$}
$a_n$ = 8($a_{n-1}$) -16($a_{n-2}$) for all n $\geq$ 2, where $a_0$ = 0, $a_1$ = 4\\

\subsubsection[f]{f - $a_n = 2(4^n) + 3n4^n$}
$a_n$ = 8($a_{n-1}$) -16($a_{n-2}$) for all n $\geq$ 2, where $a_0$ = 2, $a_1$ = 20\\

\subsubsection[g]{g - $a_n = (-4)^n$}
$a_n$ = -4($a_{n-1}$) for all n $\geq$ 1, where $a_0$ = 1\\

\subsubsection[h]{h - $a_n = n^24^n$}
$a_n$ = 12($a_{n-1}$) -48($a_{n-2}$) + 64($a_{n-3}$) for all n $\geq$ 3, where $a_0$ = 0,$a_1$ = 4, $a_2$ = 64 \\

\subsection{Question 22a - An employee joined a company in 2023 with a starting salary of \$50,000. Every year this employee receives a raise of \$1000 plus 5\% of the salary of the previous year.}
\subsubsection[a]{a - Set up a recurrence relation for the salary of this employee n years after 2023.}
$a_0$ = 50,000\\
$a_n$ = $a_{n-1}$ + 1,000 + 0.05$a_{n-1}$\\
$a_n$ = $1.05a_{n-1}$ + 1,000 for all n$\geq$1, where $a_0$ = 50,000\\

\subsection{Question 22b - What will the salary of this employee be in 2031?}
n = 2031 - 2023 = 8\\
$a_n$ = $1.05a_{n-1}$ + 1,000 for all n$\geq$1, where $a_0$ = 50,000\\
$a_1$ = $1.05a_{0}$ + 1,000 = 53,500 in 2024\\
$a_2$ = $1.05a_{1}$ + 1,000 = 57,175 in 2025\\
$a_3$ = $1.05a_{2}$ + 1,000 = 61,033.75 in 2026\\
$a_4$ = $1.05a_{3}$ + 1,000 = 65,085.44 in 2027\\
$a_5$ = $1.05a_{4}$ + 1,000 = 69,339.71 in 2028\\
$a_6$ = $1.05a_{5}$ + 1,000 = 73,806.70 in 2029\\
$a_7$ = $1.05a_{6}$ + 1,000 = 78,497.04 in 2030\\
$a_8$ = $1.05a_{7}$ + 1,000 = 83,421.89 in 2031\\
Thus, the salary in 2031 would be \$83,421.89


\subsection{Question 32 - Find the value of each of these sums.}

\subsubsection[a]{a-$\sum_{j=0}^{8} 1 + (-1)^j $}
$\sum_{j=0}^{8} 1 + (-1)^j $ = $\sum_{j=0}^{8} 1$ + $\sum_{j=0}^{8}(-1)^j $\\
$\sum_{j=0}^{8} 1 + (-1)^j $ = 9(1) + 1 + 4(-1 + 1) = 10

\subsubsection[b]{b-$\sum_{j=0}^{8} (3^j - 2^j) $}
$\sum_{j=0}^{8} (3^j - 2^j) $ = $\sum_{j=0}^{8} (3^j) $ - $\sum_{j=0}^{8} (2^j) $\\
Using the geometric series formula:\\
$\sum_{j=0}^{8} (3^j - 2^j) $ = ($\frac{3^0(1-3^9)}{1-3}$) - ($\frac{2^0(1-(2)^9)}{1-2}$)\\
$\sum_{j=0}^{8} (3^j - 2^j) $ = (9841 - 511) = 9330

\subsubsection[c]{c-$\sum_{j=0}^{8} (2 \cdot 3^j + 3\cdot 2^j) $}
Using the geometric series formula:\\
$\sum_{j=0}^{8} (2 \cdot 3^j + 3\cdot 2^j) $ = 2($\frac{3^0(1-3^9)}{1-3}$) + 3($\frac{2^0(1-(2)^9)}{1-2}$)\\
using part b, \\
$\sum_{j=0}^{8} (2 \cdot 3^j + 3\cdot 2^j) $ = 2(9841) + 3(511) = 21215

\subsubsection[d]{d-$\sum_{j=0}^{8} (2^{j+1} - 2^j) $}
$\sum_{j=0}^{8} (2^{j+1} - 2^j) $ = $\sum_{j=0}^{8} (2^{j}2 - 2^j) $\\
$\sum_{j=0}^{8} (2^{j+1} - 2^j) $ = (2$\sum_{j=0}^{8} 2^{j}) - (\sum_{j=0}^{8} (2^j)) = (\sum_{j=0}^{8} (2^j))$\\
Using the geometric series formula:\\
$\sum_{j=0}^{8} (2^{j+1} - 2^j) $ = $\frac{1-2^9}{1-2} = 511$\\


\section{Section 5.1}
\subsection{Question 8 - Prove that $2 - 2 \cdot 7 + 2 \cdot 72 - \ldots  + 2(-7)^n = (1 -(-7)^{n+1})/4$ whenever n is a nonnegative integer.}
P(n) = $2 - 2 \cdot 7 + 2 \cdot 72 - \ldots  + 2(-7)^n = (1 -(-7)^{n+1})/4$ for all n$\geq$0\\
Basis step:
$P(0) = a_0 = 2(-7)^0 = 2 = (1 -(-7)^{0+1})/4 = 2$ which holds\\
Inductive step:
$p(k) \to p(k+1)$\\
Assume p(k) holds for k $\geq$ 0, $2 - 2 \cdot 7 + 2 \cdot 72 - \ldots  + 2(-7)^k = (1 -(-7)^{k+1})/4$\\
RHS: $(1 -(-7)^{(k+1)+1})/4$\\
LHS: $(1 -(-7)^{k+1})/4 + a_{k+1}$\\
$(1 -(-7)^{k+1})/4 + 2(-7)^{k+1}$\\
$(1 -(-7)^{k+1})/4 + (8(-7)^{k+1})/4$\\
$(1 + (8-1)(-7)^{k+1})/4$\\
$(1 + 7(-7)^{k+1})/4$\\
$(1 - -7(-7)^{k+1})/4$\\
$(1 - (-7)^{(k+1)+1})/4$\\
LHS = RHS therefore $p(k) \to p(k+1)$ holds for all n $\geq$ 0 by induction.

\subsection{Question 12 - Prove $\sum_{j=0}^{n} (-\frac{1}{2})^j = \frac{2^{n+1} + (-1)^n}{3 \cdot 2^n}$ for n $\geq$ 0}
% $\sum_{j=0}^{n} (-\frac{1}{2})^j = \frac{2^{n+1} + (-1)^n}{3 \cdot 2^n}$ for n $\geq$ 0

P(n) = $\sum_{j=0}^{n} (-\frac{1}{2})^j = \frac{2^{n+1} + (-1)^n}{3 \cdot 2^n}$ for all n $\geq$ 0\\
Basis step:
$P(0) = a_0 = \sum_{j=0}^{0} (-\frac{1}{2})^0 = 1 = \frac{2^{0+1} + (-1)^0}{3 \cdot 2^0} = \frac{3}{3} = 1$ which holds\\
Inductive step:
$p(k) \to p(k+1)$\\
Assume p(k) holds for k $\geq$ 0, $\sum_{j=0}^{k} (-\frac{1}{2})^j = \frac{2^{k+1} + (-1)^k}{3 \cdot 2^k}$\\
RHS: $\frac{2^{(k+1)+1} + (-1)^{k+1}}{3 \cdot 2^{k+1}}$\\
LHS: $\sum_{j=0}^{k} (-\frac{1}{2})^j + (-\frac{1}{2})^{k+1}$\\
$\frac{2^{k+1} + (-1)^k}{3 \cdot 2^k} + (-\frac{1}{2})^{k+1}$\\
$\frac{2(2^{k+1} + (-1)^k)}{3 \cdot 2^{k+1}} + (\frac{3(-1)^{k+1}}{3\cdot2^{k+1}})$\\
$\frac{(2(2^{k+1}) -(-1)^k)}{3 \cdot 2^{k+1}}$\\
$\frac{(2^{(k+1)+1} + (-1)^{k+1})}{3 \cdot 2^{k+1}}$\\
LHS = RHS therefore $p(k) \to p(k+1)$ holds for all n $\geq$ 0 by induction.

\subsection{Question 14 - Prove that for every positive integer n, $\sum_{k = 1}^{n} (k2^k) = (n-1)2^{n+1} + 2$}
P(n) = $\sum_{k = 1}^{n} (k2^k) = (n-1)2^{n+1} + 2$ for all n $\geq 1$\\
Basis step:
$P(1) = a_1 = \sum_{k = 1}^{n} (k2^k) = 2 = (n-1)2^{n+1} + 2 = 0 + 2 $ which holds\\
Inductive step:
$p(k) \to p(k+1)$\\
Assume p(j) holds for j $\geq$ 1, $\sum_{k = 1}^{j} (k2^k) = (j-1)2^{j+1} + 2$\\
RHS: $((j+1)-1)2^{(j+1)+1} + 2$\\
$(j)2^{j+2} + 2$\\
LHS: $\sum_{k = 1}^{j} (k2^k) + (j+1)2^{j+1}$\\
$(j-1)2^{j+1} + 2 + (j+1)2^{j+1}$\\
$(2j)2^{j+1} + 2$\\
$(j)2^{j+2} + 2$\\
LHS = RHS therefore $p(k) \to p(k+1)$ holds for all n $\geq$ 1 by induction.

\subsection{Question 18 - Prove by mathematical induction: Let P(n) be the statement that $n! < n^n$, where n is an integer greater than 1}
P(n) = $n! < n^n$ for all n $>$ 1\\
a) Basis step: P(2) = $2! < 2^2$\\
b) P(2) = $2! = 2 < 2^2 = 4$ which holds\\
Inductive step:
$p(k) \to p(k+1)$\\
c) Assume p(k) holds for k $>$ 1, $k! < k^k$\\
d) Need to prove that p(k) implies p(k+1) is true, meaning the LHS < RHS for p(k+1)\\
e)
RHS: $(k+1)^{k+1}$\\
$(k+1)(k+1)^{k}$\\
LHS:$(k+1)!$\\
$(k+1)(k)!$\\
$(k+1)(k)! < (k+1)(k^k)$ Application of the hypothesis\\
$(k+1)(k)! < (k+1)(k^k)$\\
Comparing LHS and RHS:\\
$(k+1)(k)! < (k+1)(k)^k < (k+1)(k+1)^{k}$ because k $<$ k+1\\
f) These steps show under mathematical induction that starting with the smallest possible integer value for n greater than 1, (that is 2), the inequality holds. For any integer greater than or equal to 2 the inequality will continue to hold, as one integer implies the next holds the inequality.
\end{document}